% GENERATED FILE. Edit canonical lessons, then run npm run generate:books.
% canonical-source-hash: fnv1a64:7382f18d1524973d
% canonical-lessons: ZH-C315-jianjian, ZH-C315-you, ZH-W315-you-again, ZH-C315-zizai, ZH-C315-jingran, ZH-C315-rengran

\chapter{Once Again}
\label{ch:zh-once-again}
\hlchaptermodality{315}

\begin{chapteropening}
\textbf{By the end of this chapter:} \emph{I can say something happened again, that I feel at home now, and that some things are as before.}

Complete the last lesson of chapter 315: I can say something happened again, that I feel at home now, and that some things are as before.
\end{chapteropening}

\section[jiànjiàn]{jiànjiàn — gradually; little by little}

\label{lesson:ZH-C315-jianjian}



\begin{quote}
Before the new one: say the Mandarin for a cemetery, then the Mandarin for to shed tears.
\end{quote}

\subsection*{You'll want to know: jiànjiàn}
\emph{jiànjiàn} — \textquotedblleft{}gradually; little by little\textquotedblright{}.

\begin{quote}
\emph{wǒ jiànjiàn xíguàn le zhèr de shēnghuó} — Little by little I got used to life here.
\end{quote}

Its characters come later; for now it is a word for the ear.

\begin{cousinweb}
\emph{jiànjiàn} is the word \emph{jiàn} said twice, inherited from Middle Chinese \emph{dzjemX} (Guangyun, in Baxter's transcription), which glosses it 'step by step' and 'to advance'.
\end{cousinweb}

\begin{grammarlens}[title={What you've built}]
One more word for telling what happened.
\end{grammarlens}

\subsection*{Your turn}
\begin{itemize}
\raggedright
  \item \emph{Say it:} \emph{jiànjiàn}
  \item \emph{Say it:} \emph{jiànjiàn}, once more
  \item \emph{Say it:} \emph{wǒ jiànjiàn xíguàn le zhèr de shēnghuó} — Little by little I got used to life here.
  \item \emph{Recall:} say the Mandarin for a cemetery, then the Mandarin for to shed tears, then say \emph{jiànjiàn} again
\end{itemize}

\begin{tcolorbox}[breakable,colback=teal!4,colframe=teal!35!black,title={Before you move on}]
What does jiànjiàn mean? (\textquotedblleft{}Gradually; little by little\textquotedblright{}.) Say it once more.
\end{tcolorbox}
\section[\texorpdfstring{\zh{又} (yòu)}{又 (yòu)}]{\zh{又} (yòu) — yet again; it happened again}

\label{lesson:ZH-C315-you}



\begin{quote}
Before the new one: say the Mandarin for to shed tears, then the Mandarin for little by little.
\end{quote}

\subsection*{You'll want to know: \zh{又}}
\textbf{\zh{又}} — \emph{yòu} — \textquotedblleft{}yet again; it happened again\textquotedblright{}.

\begin{quote}
\textbf{\zh{你又来了}} \emph{nǐ yòu lái le} — Here you are again.
\end{quote}

The new character, in pieces: \zh{又} — two strokes: a bar that turns into a long sweep to the lower left, then a long stroke falling to the right across it, the fork at the foot of \zh{报} and \zh{发}. The next lesson puts it in your hand.

\begin{cousinweb}
\textbf{\zh{又}} (\emph{yòu}) is inherited from Middle Chinese \emph{hjuwH} (Guangyun, in Baxter's transcription), which glosses it 'further' and 'once more'.
\end{cousinweb}

\begin{grammarlens}[title={What you've built}]
One more word for telling what happened.
\end{grammarlens}

\subsection*{Your turn}
\begin{itemize}
\raggedright
  \item \emph{Say it:} \emph{yòu}
  \item \emph{Say it:} \emph{yòu}, once more
  \item \emph{Say it:} \emph{nǐ yòu lái le} — Here you are again.
  \item \emph{Recall:} say the Mandarin for to shed tears, then the Mandarin for little by little, then say \emph{yòu} again
\end{itemize}

\begin{tcolorbox}[breakable,colback=teal!4,colframe=teal!35!black,title={Before you move on}]
What does \zh{又} mean? (\textquotedblleft{}Yet again; it happened again\textquotedblright{}.) Say it once more.
\end{tcolorbox}
\section[\texorpdfstring{\zh{又} (yòu)}{又 (yòu)}]{\zh{又} — the fork at the foot of \zh{报}}

\label{lesson:ZH-W315-you-again}



\begin{quote}
The word first, then the character inside it.

\begin{itemize}
\raggedright
  \item \emph{Recall:} say \emph{yòu} — \textbf{R1}, one lesson back, before it decays
\end{itemize}

\begin{quote}
\zh{报} \zh{发}
\end{quote}

\emph{Read it:} each one aloud, with nothing said first — \emph{bào}, \emph{fā}
\end{quote}

\subsection*{Script — the fork at the foot of \zh{报}}
\hlblockfigure{\detokenize{figures/ZH-W315-you-again-filmstrip.pdf}}{How \zh{又} is written, stroke by stroke}

\begin{quote}
\zh{又}
\end{quote}

\textbf{One piece, two strokes.}

\textbf{A bar that turns into a sweep}: across to the right, then a long sweep down to the lower left.

\textbf{A long stroke falling to the right}, from under the start of the bar, across the sweep to the foot.

You have made this fork in \zh{报}, \zh{发} and \zh{经}.

\subsection*{Your turn}
\begin{itemize}
\raggedright
  \item \emph{Write it:} \zh{又}, all two, with the model in front of you
  \item \emph{Cover:} write it again from memory, and count the pieces, not the strokes
\end{itemize}

\begin{tcolorbox}[breakable,colback=teal!4,colframe=teal!35!black,title={Before you move on}]
How many strokes in \zh{又}? (\textbf{Two}.) Where have you made this shape? (\textbf{At the foot of \zh{报}, and in \zh{发} and \zh{经}}.)

Source: \href{https://raw.githubusercontent.com/chanind/hanzi-writer-data/68d10a4b21150cae5e1ebbd223eed289cf32d90c/data/\%E5\%8F\%88.json}{Hanzi Writer Data \zh{又}, pinned stroke paths and medians}.
\end{tcolorbox}
\section[\texorpdfstring{\zh{自在} (zìzai)}{自在 (zìzai)}]{\zh{自在} (zìzai) — at ease; free and relaxed}

\label{lesson:ZH-C315-zizai}



\begin{quote}
Before the new one: say the Mandarin for little by little, then the Mandarin for it happened again.
\end{quote}

\subsection*{You'll want to know: \zh{自在}}
\textbf{\zh{自在}} — \emph{zìzai} — \textquotedblleft{}at ease; free and relaxed\textquotedblright{}.

\begin{quote}
\textbf{\zh{我来了一年才觉得自在}} \emph{wǒ lái le yì nián cái juéde zìzai} — It took me a year to feel at home.
\end{quote}

\begin{cousinweb}
\textbf{\zh{自在}} (\emph{zìzai}) is built from \zh{自} 'by oneself' and \zh{在} 'to be in place', both inherited from Middle Chinese, \emph{dzijH} and \emph{dzojX} (Guangyun, in Baxter's transcription), which glosses the first 'from' and 'by way of'; the sense 'oneself' is not given there.
\end{cousinweb}

\begin{grammarlens}[title={What you've built}]
One more word for telling what happened.
\end{grammarlens}

\subsection*{Your turn}
\begin{itemize}
\raggedright
  \item \emph{Say it:} \emph{zìzai}
  \item \emph{Say it:} \emph{zìzai}, once more
  \item \emph{Say it:} \emph{wǒ lái le yì nián cái juéde zìzai} — It took me a year to feel at home.
  \item \emph{Recall:} say the Mandarin for little by little, then the Mandarin for it happened again, then say \emph{zìzai} again
\end{itemize}

\begin{tcolorbox}[breakable,colback=teal!4,colframe=teal!35!black,title={Before you move on}]
What does \zh{自在} mean? (\textquotedblleft{}At ease; free and relaxed\textquotedblright{}.) Say it once more.
\end{tcolorbox}
\section[jìngrán]{jìngrán — to my surprise; unexpectedly}

\label{lesson:ZH-C315-jingran}



\begin{quote}
Before the new one: say the Mandarin for it happened again, then the Mandarin for at ease in a place.
\end{quote}

\subsection*{You'll want to know: jìngrán}
\emph{jìngrán} — \textquotedblleft{}to my surprise; unexpectedly\textquotedblright{}.

\begin{quote}
\emph{tā jìngrán wàng le wǒ de shēngrì} — He actually forgot my birthday!
\end{quote}

Its characters come later; for now it is a word for the ear.

\begin{cousinweb}
\emph{jìngrán} is built from \emph{jìng} 'to the end' and \zh{然} 'so', both inherited from Middle Chinese, \emph{kjaengH} and \emph{nyen} (Guangyun, in Baxter's transcription), which glosses the first 'to use up' and 'the end'; how the pair came to mean 'to one's surprise' is not settled.
\end{cousinweb}

\begin{grammarlens}[title={What you've built}]
One more word for telling what happened.
\end{grammarlens}

\subsection*{Your turn}
\begin{itemize}
\raggedright
  \item \emph{Say it:} \emph{jìngrán}
  \item \emph{Say it:} \emph{jìngrán}, once more
  \item \emph{Say it:} \emph{tā jìngrán wàng le wǒ de shēngrì} — He actually forgot my birthday!
  \item \emph{Recall:} say the Mandarin for it happened again, then the Mandarin for at ease in a place, then say \emph{jìngrán} again
\end{itemize}

\begin{tcolorbox}[breakable,colback=teal!4,colframe=teal!35!black,title={Before you move on}]
What does jìngrán mean? (\textquotedblleft{}To my surprise; unexpectedly\textquotedblright{}.) Say it once more.
\end{tcolorbox}
\section[réngrán]{réngrán — as before; even now}

\label{lesson:ZH-C315-rengran}



\begin{quote}
Before the new one: say the Mandarin for at ease in a place, then the Mandarin for to my surprise.
\end{quote}

\subsection*{You'll want to know: réngrán}
\emph{réngrán} — \textquotedblleft{}as before; even now\textquotedblright{}.

\begin{quote}
\emph{shí nián yǐhòu, tā réngrán zhù zài zhèr} — Ten years later, he still lives here.
\end{quote}

Its characters come later; for now it is a word for the ear.

\begin{cousinweb}
\emph{réngrán} is built from \emph{réng} 'to follow on' and \zh{然} 'so', both inherited from Middle Chinese, \emph{nying} and \emph{nyen} (Guangyun, in Baxter's transcription), which glosses the first 'to follow', 'to repeat' and 'often'.
\end{cousinweb}

\begin{grammarlens}[title={What you've built}]
One more word for telling what happened.
\end{grammarlens}

\subsection*{Your turn}
\begin{itemize}
\raggedright
  \item \emph{Say it:} \emph{réngrán}
  \item \emph{Say it:} \emph{réngrán}, once more
  \item \emph{Say it:} \emph{shí nián yǐhòu, tā réngrán zhù zài zhèr} — Ten years later, he still lives here.
  \item \emph{Recall:} say the Mandarin for at ease in a place, then the Mandarin for to my surprise, then say \emph{réngrán} again
\end{itemize}

\begin{tcolorbox}[breakable,colback=teal!4,colframe=teal!35!black,title={Before you move on}]
What does réngrán mean? (\textquotedblleft{}As before; even now\textquotedblright{}.) Say it once more.
\end{tcolorbox}
